\section[nosectionpage]{How to Create a Webserver Object}

\begin{frame}[fragile]{How to Create a Webserver Object}
  After including a concrete \inlinecppcode|"express/[security]/[network]/WebApp.h"| header one can instanciate a webserver object by using the class name \inlinecppcode|express::[security]::[network]::WebApp| and either the \textcolor{fh-ooe-red}{default constructor} or the \textcolor{fh-ooe-red}{constructor expecting} an \textcolor{fh-ooe-red}{instance name} as argument.
  The \textcolor{fh-ooe-red}{instance name} must be \textcolor{fh-ooe-red}{unique} in one server application. Thus, you can \textcolor{fh-ooe-red}{not} create \textcolor{fh-ooe-red}{two WebApp instances} with the \textcolor{fh-ooe-red}{same instance name}.

  \begin{cppcode}[gobble=4]
    express::[security]::[network]::WebApp app;             // Anonymous WebApp
    express::[security]::[network]::WebApp app("mywebapp"); // Named WebApp
  \end{cppcode}
  \begin{description}
    \item [Anonymous WebApp:] No command line arguments are provided.
    \item [Named WebApp:] Command line arguments are provided automatically by SNode.C.
  \end{description}
  The \textcolor{fh-ooe-red}{API} provided by any of the \textcolor{fh-ooe-red}{WebApp} classes is the same and very \textcolor{fh-ooe-red}{similar} to the \textcolor{fh-ooe-red}{API of NodeJS-Express}. You can use e.g.
  \begin{itemize}
    \item \inlinecppcode|app.use(...)|
    \item \inlinecppcode|app.get(...)|
    \item \inlinecppcode|app.post(...)|
    \item \inlinecppcode|app.put(...)|
  \end{itemize}
  as in NodeJS-Express but with C++ syntax.
\end{frame}
